Cybersickness remains a major barrier to the widespread adoption of virtual reality (VR).
Current subjective questionnaires are retrospective and insufficient for the dynamic, continuous assessment required to diagnose its determinants and implement real-time mitigation.

To address this, the field has increasingly investigated neuromarkers: signals derived from non-invasive neuroimaging as well as neurostimulation.
This PRISMA-guided systematic review synthesizes 92 studies investigating the neurophysiology of cybersickness in VR.
The corpus explicitly unites recording modalities (77 electroencephalography, 4 functional near-infrared spectroscopy) with neurostimulation interventions (11 studies spanning galvanic vestibular, transcranial direct-current, and transcranial alternating-current stimulation) under a unified frame.

Synthesis of the literature reveals consistent neural correlates alongside a critical methodological confound.
Cybersickness is predominantly associated with increased low-frequency (delta and theta) power and decreased alpha-band power.
However, we report meta-analytic evidence that VR hardware significantly modulates these band-power responses ($p=0.004$ on alpha, large effect; $p=0.015$ on beta), indicating that prior contradictions in the literature are partly attributable to hardware confounds rather than to differences in cybersickness per se.

Despite these advances, the literature faces substantial methodological and reporting barriers.
A methodological audit of the 38 classification or prediction studies reveals sparse reporting and arbitrary thresholds: only one study reports leave-one-subject-out validation, nineteen do not disclose a sick-versus-non-sick decision threshold, and thirteen report accuracy without sensitivity, specificity, AUC, or F1.
We conclude with a call for standardized reporting, multi-site open data, and causal validation via interventional designs to establish robust neuromarkers that can enable adaptive, user-aware immersive systems.
